\section{The exponent from the OpenAI parameters as displayed}\label{app:openai-constant}

\paragraph{What this appendix proves and how to read it.}
Theorem~1.1 of \cite{OpenAI} asserts $|A|\le CN^{1-c}$ for an
unspecified absolute $c>0$ (their $c$ is our $\eps$). This appendix
evaluates $c$ for \emph{one displayed choice} of their parameters:
the fixed construction of their \S3.1, the threshold $P$
beyond which their conditional-mixing estimate (a bound on
conditional correlations given the other labels) applies to the
prime cutoff, and their stated lower bound $\rho\ge M_0^{-23}$
for the proportion $\rho$ of residue classes modulo $M_0$ that
are good (admit the required square increments, their
Lemma~8.1). Equation
numbers such as (3.16) and (9.13) refer to \cite{OpenAI}, not to
this paper. Two scope warnings are essential. First, this is not a
claim about the best exponent obtainable from their method with
different parameters or sharper local estimates; it quantifies only
the choices as displayed. Second, the integers $d$ and $e_0$ below
are exact, but the prime product in $M_0$ is left unevaluated and
logarithms are asymptotic approximations, not interval
certificates. Nothing in Parts~I--V depends on this appendix; it
justifies only the comparison in the abstract and
\cref{subsec:intro:new}. The computation is reproduced by a short
standard-library program in the repository accompanying this paper
(\path{analysis/openai_constant/compute_openai_constant.py}); rerunning
it reproduces
\path{analysis/openai_constant/openai_constant_results.json}
byte for byte.

\paragraph{Result.}
For this specialization the exponent is, symbolically,
\begin{equation}\label{eq:oa:symbolic}
  \eps_{\mathrm{OA}}
   =\frac1d\cdot
     \frac{-\log\bigl(1-\tfrac34 M_0^{-23}\bigr)}{4\log M_0},
   \qquad
   M_0=8\!\!\prod_{3\le p<P}\!p,\qquad
   P=\big\lceil(C_{\rm lead}+C_{\rm mix})^{32}\big\rceil,
\end{equation}
with $d=(24!)^{161}$ and $C_{\rm lead},C_{\rm mix}$ as in
\eqref{eq:oa:clead} below. At the level of the asymptotic estimates
proved below,
\begin{equation}\label{eq:oa:tower}
  \eps_{\mathrm{OA}}=2^{-2^{X}},\qquad
  X=96e_0+16\log_2e_0+\log_2\frac{23}{\ln2}+176+o(1)
   =96e_0+307462.2929\ldots,
\end{equation}
where
\begin{equation}\label{eq:oa:e0}
  e_0=(24!)^{161}\sum_{j=81}^{161}\binom{161}{j}(24!-1)^{161-j}
     =2.01561769398235743936\ldots\times10^{5781}
\end{equation}
is the exact number of directed square constraints (edges) in
their graph $G_\emptyset$ (their `leading list graph': vertices
the $d$ tuples of orderings, with an edge $v\to w$ when $w$ agrees
with the prescribed translate of $v$ in at least $81$ of the $161$
blocks; \cref{subsec:oa:fixed}); $e_0$ has $5782$ decimal
digits. In \eqref{eq:oa:tower} the $o(1)$ term, and the
$O(\cdot)$ terms below, are as the corresponding parameter
($e_0$, $P$ or $M_0$) tends to infinity in the estimates of
\cite{OpenAI}; at the displayed finite values they are not
enclosed (second scope warning). At that asymptotic level
$\log_{10}\log_{10}\log_{10}(1/\eps_{\mathrm{OA}})$ is
$5782.7652891675\ldots$ (digits of the asymptotic expression,
not certified digits of the exact defined value), so
\begin{equation}\label{eq:oa:decimal}
  \eps_{\mathrm{OA}}\approx10^{-10^{10^{5782.7653}}},
\end{equation}
a triple exponential. The decimals in the top exponent matter:
replacing $5782.7653$ by $5782$ changes the value by many orders of
magnitude on a double-logarithmic scale, so \eqref{eq:oa:decimal}
should always be quoted with them, or as
$\log_{10}\log_{10}\log_{10}(1/\eps_{\mathrm{OA}})\approx5782.77$.
Using the exact good-residue proportion $\rho$ in place of
$\rho\ge M_0^{-23}$ changes only the additive constant in $X$ (to
$96e_0+288250.69\ldots$); the tower height and leading term $96e_0$
are unchanged.

\subsection{The fixed construction}\label{subsec:oa:fixed}

Their \S3.1 and (3.1) fix $h=24$, $t_0=80$ and
$\ell=2t_0+1=161$ blocks. Write $K=24!=620448401733239439360000$
for the number of orderings in one block, $d=K^\ell=(24!)^{161}$
($\log_{10}d=3830.6256\ldots$) for the number of their points
(their analogue of slots; their $161$ blocks are blocks of the
ordering construction, not prime blocks of this paper),
$r=K/2$, and $\delta=1/(10h^2)=1/5760$. The remaining displayed
quantities are the marked-cut bound
$\kappa=\tfrac9{10}\delta^{h-1}\approx2.9\cdot10^{-87}$ (their
Lemma~3.3; $\kappa$ is used below only through the constant $c_*$
in $P_{\rm ref}$), the index-weight bound
$B=\exp\bigl(d(t_0+1)\binom h2\delta/2\bigr)$ with
$\log_2B=1.1823\cdot10^{3831}$ (their (3.8)), and the out-degree
$\Delta=\sum_{j\ge81}\binom{161}{j}(K-1)^{161-j}$
($\log_{10}\Delta=1950.6788\ldots$) in $G_\emptyset$, so that
$e_0=d\Delta$ as in \eqref{eq:oa:e0}. The count for $\Delta$ is
exact: in each block a vertex $w$ either equals the one
prescribed translate of $v$ in their definition of
$G_\emptyset$ (their \S3.1; one choice) or not ($K-1$ choices),
and an edge
$v\to w$ requires at least $t_0+1=81$ matching blocks.

\subsection{The prime threshold $P$}\label{subsec:oa:P}

Their proof takes one integer $P$ exceeding every threshold
(Proposition~3.1, Lemma~3.6, (7.1) of \cite{OpenAI}). In
$\log_2$ size the candidates are: $P_{\rm ref}$, the threshold
of their (3.16) (its constants $C_*,c_*,n$ are those of their
Proposition~3.1; the formula is not needed here), for which the
program computes $\log_2P_{\rm ref}=2.26\cdot10^{3835}$ (well
below $2^{e_0}$ on a logarithmic scale in the sense checked by the
program); $P_{\rm mass}=(2^{e_0+1}e_0)^2$ with
$\log_2P_{\rm mass}=2e_0+38412.2$ (their (3.19));
$C_{\rm exc}^{16/17}$ with $C_{\rm exc}=c_0^{-1}\binom{161}{81}B$,
$c_0=2^{-e_0-1}$ (their (3.22)), for which the program computes
$\log_2C_{\rm exc}=e_0+1.1823\cdot10^{3831}$; the Lemma~3.6 threshold
$(C_{\rm lead}+C_{\rm mix})^{32}$; and the Gauss-bound threshold of
their (7.1), with $\log_2$ below $7$. The Lemma~3.6 threshold is
the largest of these: the program checks $\log_2P_{\rm ref}<e_0$
and $\log_2C_{\rm exc}<2e_0$, and with
$\log_2P_{\rm mass}=2e_0+38412.2$ each candidate is below the
value $96e_0+176+16\log_2e_0$ of \eqref{eq:oa:logP} below. There
\begin{equation}\label{eq:oa:clead}
  C_{\rm lead}=c_0^{-1}\bigl(2^{e_0+1}+2^{2e_0+2}\bigr)\sqrt{32e_0},
  \qquad C_{\rm mix}=2^\ell B/c_0,
\end{equation}
so $\log_2C_{\rm lead}=3e_0+5.5+\tfrac12\log_2e_0+O(2^{-e_0})$ and
$C_{\rm mix}/C_{\rm lead}<2^{-e_0}$. The requirement
$C_{\rm cond}p^{-1/16}\le p^{-1/32}$ for $p\ge P$, with
$C_{\rm cond}=C_{\rm lead}+C_{\rm mix}$, gives
$P=\lceil C_{\rm cond}^{32}\rceil$, i.e.
\begin{equation}\label{eq:oa:logP}
  \log_2P=96e_0+176+16\log_2e_0+O(2^{-e_0}).
\end{equation}
The mechanism is worth stating in words: the leading component of
their tuple law has mass about $2^{-e_0}$, and the
conditional-mixing estimate divides by that mass. Even with the
exact degree in place of the crude bound used there, this division
alone forces $\log_2P\ge32e_0$ at this step. A different mixing
estimate need not contain the factor $c_0^{-1}=2^{e_0+1}$;
removing this normalisation
is one of the structural changes in the present paper
(\cref{subsec:intro:new}).

\subsection{Small modulus, good residues, and the exponent}\label{subsec:oa:exponent}

Write $M_0=8\prod_{3\le p<P}p$, so
$\ln M_0=\ln4+\vartheta(P-1)=P\bigl(1+O((\log P)^{-3})\bigr)$ by
explicit Chebyshev bounds, where $\vartheta$ is Chebyshev's
function. Their Lemma~8.1 gives $\rho\ge M_0^{-23}$ for the
good-residue proportion. With their parameters
$\beta=10^{-5}$, $\gamma=1/64$,
$\tau=\beta\gamma/(1000(d+1))$ and $\sigma_0=\gamma\tau/4$
(their (6.1), (7.3); these letters are OpenAI's, not parameters
of this paper), their (9.13) sets their
$a_0$ (not the evaluation exponent $a_0$ of this paper) to
$\min\{\tfrac14,\sigma_0/2,\log(1/r_{\rm OA})/(4\log M_0)\}$
with $r_{\rm OA}=1-\tfrac34\rho$ (not the $r=K/2$ above), and
then $\eps_{\mathrm{OA}}=a_0/d$. The
analytic entry alone would give
$\sigma_0/(2d)=1.711\ldots\cdot10^{-7674}$, enormous compared with
the residue entry, so the third entry is the minimum. Hence
\[
  \log_2\frac1{\eps_{\mathrm{OA}}}
   =23\log_2M_0+\log_2d+\log_2\frac{16}3+\log_2\ln M_0
     +O(M_0^{-23}),
\]
and hence
\[
 X=\log_2\log_2(1/\eps_{\mathrm{OA}})
 =\log_2P+\log_2(23/\ln2)
   +O\!\left((\log P)^{-3}+\frac{\log P+\log d}{P}\right).
\]
The first error term retains the relative error in $\ln M_0$;
the second comes from the additive logarithmic terms before taking
the outer logarithm. Together with \eqref{eq:oa:logP}, these are
$o(1)$ in \eqref{eq:oa:tower}. Converting bases gives
$\log_{10}\log_{10}\log_{10}(1/\eps_{\mathrm{OA}})
=5782.7652891675\ldots$, i.e.\ \eqref{eq:oa:decimal}; at the
displayed values the omitted asymptotic corrections change this
triple logarithm by about $10^{-5779}$ (the program's estimate,
not an interval enclosure), so the quoted $5782.77$ is
unaffected.

For comparison, the exact proportion is
$\rho=8^{-23}\prod_{3\le p<P}R_{24}(p)$ with
\[
  R_{24}(p)=
  \frac{(p-1)^{24}+(p-1)\sum_{j=0}^{12}\binom{24}{2j}
        (\epsilon_pp)^j}{(2p)^{24}},
  \qquad\epsilon_p=(-1)^{(p-1)/2},
\]
the cycle-closing probability for $24$ uniform residue labels with
nonzero-square increments. This gives
$\log_2(1/\rho)=24\pi(P)+O(\log\log P)$ and therefore
$X=96e_0+288250.69\ldots$ in place of \eqref{eq:oa:tower}: the
same tower with a different additive constant, invisible in
\eqref{eq:oa:decimal} at the displayed precision.

The implied constants in their Propositions~6.1, 7.2, 8.2 and~9.1
affect only the prefactor $C$ and the starting scale; they never
enter $a_0$, which is fixed before the induction. They therefore
leave no unspecified parameter in \eqref{eq:oa:symbolic}. This
does not make \eqref{eq:oa:decimal} an exact numerical evaluation
(the product defining $M_0$ has about $P/\log P$ factors and is
never formed), nor an optimal choice within their framework.

\paragraph{Consequences.} Nothing in Parts~I--V uses this
appendix. Its chain $e_0\to c_0\to C_{\rm cond}\to P\to M_0\to
\rho\to a_0$ shows where the tower comes from: not from the degree
$d$ alone, but from dividing by the mass $c_0$ of the leading
graph in the mixing step, which forces $\log_2P\ge32e_0$, and
then dividing by $\log M_0$, of size $P$, in their $a_0$. The present paper
removes the first normalisation by relative (density-normalised)
estimates and makes the factor coming from good residues
independent of $M_0$ by coupling; these two changes remove the
division by $c_0$ in $P$ and the division by $\log M_0$ in their
$a_0$, i.e.\ two of the three exponentials in
\eqref{eq:oa:decimal}, before any local constant is optimised;
see \cref{subsec:intro:new}.
